\documentclass[12pt]{article}

\usepackage[T1]{fontenc}
\usepackage[utf8]{inputenc}
\usepackage[english]{babel}
\usepackage{amsmath,amssymb,amsthm}
\usepackage[margin=1.25in]{geometry}
\usepackage{setspace}
\usepackage[authoryear,round]{natbib}
\usepackage[colorlinks=true,linkcolor=black,citecolor=blue,urlcolor=blue]{hyperref}

\onehalfspacing
\setlength{\parindent}{1.5em}
\setlength{\parskip}{0pt}

\newtheorem{proposition}{Proposition}[section]

\title{Common Decency and the Segmented Self:\\
Liberal Doubling, Reciprocity, and the Limits of a Conservative Left}
\author{Flyxion\\Independent Researcher}
\date{October 2026}

\begin{document}
\maketitle

\begin{abstract}
\noindent Jean-Claude Michéa argues that economic liberalism and cultural liberalism are two expressions of a single philosophical project, and that the second prepares individuals for the first by dissolving the inherited obligations that would otherwise resist the extension of the price form. The argument has recently reached a wide English-language audience through a long, informal video lecture whose treatment of Michéa mixes his anthropological claims with the lecturer's own polemical vocabulary. This essay separates those layers and examines each in turn, drawing on sociology, moral and social psychology, political theory, and literary theory. It distinguishes three propositions that the doubling thesis tends to run together: that capitalism can accommodate cultural reform, that it can commodify the identities and relationships associated with reform, and that cultural reform is causally necessary to capitalist expansion. The first two are well supported. The third is not established by the lecture, and the symmetry of the market's treatment of tradition and emancipation weighs against it, although that symmetry leaves open historically specific causal relations, which require case-by-case study. What survives is a general claim about representations that come to substitute for the objects they describe, a mechanism that links Orwell's prefabricated phrases, Hacking's looping classifications, and the market projection of relationships onto prices. The essay states a formal criterion of viable continuation, with an explicit time horizon, a threshold of viability, a scope covering every affected agent and not only members, and exit and voice built into the definition, and argues that the criterion supplies the element missing from Michéa's conservatism. The essay does not claim to prove the doubling thesis, and it does not claim to refute it.
\end{abstract}

\section{Introduction}

A long lecture on a little-known French philosopher, delivered in an informal register by a speaker who admits uncertainty and digresses freely, has been watched more than a hundred thousand times. That fact is minor, but it frames the question the lecture raises. The lecture presents Jean-Claude Michéa as a thinker who left the French Communist Party, spent his career teaching philosophy to adolescents rather than university students, retired to a village, declined television and the commentary economy, and wrote for decades that the political left has mistaken the dissolution of limits for the enlargement of freedom. The question raised is whether this thesis, in the form in which it circulates, is a contribution to the understanding of how market societies reproduce themselves or an instrument of cultural antagonism that borrows the authority of social theory.

The question is difficult because the thesis is at once empirical, normative, and interpretive. It is empirical in claiming that certain cultural transformations have causal relations to market expansion. It is normative in claiming that the relationships dissolved by those transformations were valuable. It is interpretive in claiming that ordinary moral life contains a resource, called common decency after Orwell, which political movements either draw upon or betray. Each of these claims belongs to a different discipline, and each is usually defended with the methods of another. The empirical claim is defended through cultural criticism, the normative claim through anecdote, and the interpretive claim through an opposition between ordinary people and intellectual elites that is itself in need of empirical support.

The essay proceeds by treating the lecture as a document of reception. Section~\ref{sec:source} separates Michéa's claims from the lecturer's additions. Section~\ref{sec:doubling} reconstructs the doubling thesis and tests it against three propositions and one qualification. Sections~\ref{sec:decency} and~\ref{sec:gift} examine common decency and the opposition between gift and contract, using the sociological literature on moral economy, embeddedness, and the commodification of intimacy. Section~\ref{sec:labels} turns to the psychology of identity labels, and Section~\ref{sec:literary} to the literary-theoretical dimension, which is more central to Michéa's lineage than the lecture's emphasis on politics suggests. Section~\ref{sec:mechanism} argues that the phrase, the label, and the price are instances of a single mechanism, in which a representation substitutes for the object it describes. Section~\ref{sec:politics} addresses the political claim that class has been displaced by recognition, and Section~\ref{sec:criterion} states the criterion of viable continuation in a form that fixes its horizon, its scope, and its treatment of exit and voice. The conclusion draws the results together.

\section{Source, Attribution, and Method}
\label{sec:source}

The lecture, titled \emph{Why Consumerism Destroyed the Left} and published by the channel Academic Agent on 8 October 2026 \citep{academicagent}, is the final installment in a weekly series on supposedly hidden conservatism among left-wing thinkers. Orwell, Pasolini, and Mark Fisher precede Michéa in that series. The lecturer works from an English-language selection of Michéa's writings and from his own long-running concepts, and the two are not always distinguishable in the delivery.

Some elements are plainly the lecturer's own. These include the notion of a ``Boomer Truth Regime,'' a description of contemporary identity politics as audience segmentation lived as a program, the comparison of political labels to membership of a pop group, and the reading of an animated film as an allegory of corporate absorption. Other elements are attributed to Michéa and have recognizable sources in his published work: the distinction between two liberalisms, the claim that the left has substituted a theology of progress for the social question, the Orwellian notion of common decency, the Maussian opposition between gift and contract, and the diagnosis of a politics organized around the lesser evil, which is the subject of \emph{L'Empire du moindre m\'al} \citep{michea}. The lecturer also reports that Michéa has declined to denounce those on the political right who cite him, which bears on the politics of reception but not on the philosophical claims.

The method adopted here is conservative about attribution. Where the lecture presents a claim as Michéa's, the claim is treated as his provisionally. Where it presents a claim in the lecturer's voice, it is treated as a contribution of the lecture and evaluated on its own. Neither is assumed to be true. The essay also declines to infer anything about the audience from the view count alone, since a count of views cannot distinguish interest in philosophical exposition from interest in political confirmation or from the pull of a confrontational title.

\section{The Doubling Thesis and Its Three Propositions}
\label{sec:doubling}

Michéa's central claim, as the lecture reports it, is that economic liberalism and cultural liberalism are not opposed. Economic liberalism privileges exchange, private ownership, and the extension of the market into every domain. Cultural liberalism privileges autonomy, self-definition, and the removal of limits that precede and exceed individual choice. The claim is that these are two faces of one project, the liberation of the individual from externally imposed constraint, and that the second prepares the individual to be the subject of the first. A relationship sustained by kinship or neighborhood reciprocity becomes a purchased service once its obligations have been declared illegitimate. An identity situated in a durable community becomes a collection of preferences once the community has lost its authority to define it.

The thesis is structurally a version of what Albert Hirschman called the self-destruction thesis in his survey of rival interpretations of market society \citep{hirschman82}. On the doux commerce thesis, markets civilize. On the feudal-remnant thesis, markets depend on virtues they did not create. On the self-destruction thesis, markets erode the moral foundations on which they depend. Marx, Polanyi, Schumpeter, and Durkheim in different ways held the third view, and Michéa belongs to that lineage. What is distinctive in his formulation is the claim that a particular cultural program, and not only the market itself, is the solvent.

The thesis is therefore best assessed by separating three propositions that the lecture tends to treat as one.

\begin{proposition}[Accommodation]
Capitalist institutions can accommodate cultural reforms associated with the removal of inherited restrictions, and often do so without loss of profitability.
\end{proposition}

\begin{proposition}[Commodification]
Capitalist institutions can commodify the identities, relationships, and symbols associated with those reforms, converting them into instruments of product differentiation, audience targeting, and brand positioning.
\end{proposition}

\begin{proposition}[Necessity]
Cultural reforms of this kind are causally necessary for capitalist expansion, in the sense that the expansion could not proceed, or would proceed substantially less, in their absence.
\end{proposition}

The first two propositions are well supported. \citet{boltanski} documented in detail how the artistic critique of the 1960s and 1970s, with its demands for authenticity, autonomy, and release from hierarchy, was absorbed into managerial discourse and reorganized as a new spirit of capitalism, in which flexibility, creativity, and personal fulfilment became sources of value. \citet{fraser2017} described an alliance between mainstream emancipatory movements and high-end service and finance sectors, which she called progressive neoliberalism, and argued that it had given recognition politics a form compatible with the dismantling of industrial welfare protections. These accounts support accommodation and commodification. They also supply a caution against the third proposition, because both present the absorption as a historical outcome of contingent alliances, not as a logical consequence of the reforms themselves.

Necessity is the weakest link, and the history of social theory contains a competing hypothesis. \citet{bell} argued that the cultural contradictions of capitalism ran in the opposite direction: a hedonistic, antinomian culture generated by consumer abundance undermines the disciplined work ethic that production requires. On that account cultural individualism is not the servant of the market but a dangerous product of it. Whether Bell or Michéa is closer to the truth is an empirical question that neither settles, and the existence of the disagreement shows that compatibility between a cultural form and market expansion does not establish dependence of the latter on the former.

There is also an observation about symmetry that bears on the third proposition. Commercial capitalism sells tradition as readily as it sells emancipation. Weddings, religious festivals, national symbolism, heritage tourism, family-centered advertising, and nostalgia form large and durable markets. If commodification were a function of the liberationist character of cultural content, traditional content would resist it, and it does not. The more defensible conclusion is that markets are indifferent to the cultural content of what they mediate, provided it can be mediated by a price, and this weighs against a universal claim that emancipatory reform is uniquely suited to market expansion. This is a modification of Michéa's thesis, and a significant one. It moves the explanatory weight from the specific content of cultural liberalism to a general process: the conversion of socially situated relationships into representations that can be priced, targeted, and exchanged.

The symmetry argument has a limited reach, and the limit should be stated. It tells against a universal necessity claim, on which every emancipatory reform is a condition of market expansion. It does not tell against historically specific causal relations, on which a particular reform facilitated a particular form of market expansion in a particular setting. \citet{fraser2016} argued that capitalist accumulation both depends on and destabilizes the unpaid work of social reproduction, and that the entry of women into paid employment, which feminist movements pursued for reasons of their own, coincided with the replacement of the family wage by the two-earner household and with the growth of marketed care. Whether the reform facilitated that expansion, was exploited by it, or merely coincided with it is the kind of question that case-specific historical work can answer and that a universal claim cannot. The essay accordingly distinguishes two claims that the third proposition could be taken to assert. The universal claim is rejected here for the reasons given. The specific claims are neither rejected nor asserted, since each depends on historical evidence about a particular reform, a particular market, and a particular period, and the symmetry of the market's appetite for tradition is not evidence about any of them.

\section{Common Decency: Moral Psychology and Its Sociological Conditions}
\label{sec:decency}

The second major component is the concept of common decency, taken from Orwell. In the lecture it names a moral floor in ordinary life: loyalty, keeping one's word, refusal to swindle a neighbor, and shame at cruelty. Orwell's own treatment of the English working class \citep{orwell37,orwell41} combined sympathetic observation with strong claims about habits of fairness and restraint, and Michéa extends the point into a general thesis that socialism drew its moral authority from such habits and cannot survive their disappearance. The summary line attributed to him in the lecture is that socialism will be decent or it will not be.

Three questions arise. The first is whether such a moral floor exists in a form that bears the weight placed on it. The moral psychology of \citet{haidt} offers partial corroboration, since loyalty, fairness, and care are among the intuitions he identifies as widely distributed across moral cultures. But the correspondence is loose, and Haidt's work is itself contested. Social psychology supports the existence of strong reciprocity norms in many settings, but it equally supports the existence of strong ingroup favoritism, which is a feature of the same moral equipment. The floor is a floor with a variable ceiling.

The second question concerns the historical sociology. \citeauthor{thompson}'s (\citeyear{thompson}) study of the moral economy of the English crowd showed that popular protest in the eighteenth century often invoked customary norms of fair price and fair dealing against market practices that violated them. This is the clearest scholarly support for the view that ordinary people held, and acted on, normative expectations that markets overrode. \citeauthor{granovetter}'s (\citeyear{granovetter}) account of embeddedness likewise shows that economic action is mediated by networks of trust and obligation that standard models treat as given. These contributions give common decency an empirical grounding that the lecture's anecdotes do not.

The third question concerns what can be inferred from the existence of such norms. A community that cultivates reciprocity can enforce exclusion and conformity alongside it. Dense solidarity among recognized members often accompanies hostility toward outsiders, and obligations can be distributed unequally, as when unpaid caregiving is assigned to particular family members within a highly reciprocal community. This means that social cohesion and moral legitimacy are not equivalent, and Michéa's defense of limits requires a criterion that distinguishes limits worth preserving from limits that merely entrench domination.

Hirschman's earlier work on exit, voice, and loyalty \citep{hirschman70} suggests a candidate. A constraint can be assessed by whether those subject to it retain meaningful avenues of voice within the arrangement and meaningful possibilities of exit from it. An obligation that sustains the agency of its participants leaves them able to contest its terms and, in the last resort, to leave. An obligation that makes participation conditional on submission closes both. On this test, common decency is the set of reciprocal expectations that survive the availability of exit and voice. Practices that persist only because exit and voice are foreclosed belong to a different category, whatever their antiquity. Section~\ref{sec:criterion} builds this test into a formal definition.

There is a further sociological precedent for the defense of limits that the lecture does not mention. \citet{durkheim} argued that desires without regulation do not liberate but produce anomie, because satisfaction depends on a sense of measure that only social norms supply. Michéa's claim that a limit is not a memory, and that limits can protect rather than oppress, is in this respect Durkheimian. It does not follow that every limit is protective, and Durkheim himself distinguished anomic from fatalistic regulation, the latter being excessive and oppressive. The Durkheimian framework therefore supports the form of Michéa's argument while reinforcing the need for a criterion.

A final difficulty concerns the opposition between ordinary people and the intellectual class. The lecture repeatedly presents a morally grounded populace and a morally detached intelligentsia, with the latter having broken faith with the former. The categories are underspecified. Education, occupation, geography, and income do not map cleanly onto moral character, and the selection of admirable working-class practices and objectionable elite practices builds the conclusion into the examples. A rigorous version of the thesis would study the actual distribution of reciprocal behavior and institutional obligation across social positions. The opposition may describe a real divergence in some settings, and the present essay does not deny it, but it is a hypothesis that the lecture asserts rather than tests.

\section{Gift and Contract}
\label{sec:gift}

Michéa's use of \citet{mauss} is the most philosophically productive element in the lecture. Mauss argued that gift exchange in many societies creates and sustains relationships through obligations to give, to receive, and to reciprocate, and that these obligations are total social facts that bind economic, legal, moral, and religious dimensions at once. The contrast with contract is that a contract specifies terms and ordinarily discharges the obligation on completion, while a gift opens a relationship whose significance depends on its continuation.

The lecturer states the opposition forcefully, replacing the abstract social contract with a social gift that binds real people locally. The opposition is analytically valuable and sociologically overdrawn. \citet{zelizer} analyzed the many ways in which intimate relationships and monetary transactions interpenetrate, rejecting both the hostile worlds view, on which money corrupts intimacy, and the view that intimacy is merely another market. People routinely combine payment and care, and the combination is managed through what she calls differentiated ties, with distinct meanings attached to distinct kinds of payment. A contract can protect a vulnerable person from arbitrary obligation, and a gift can create a debt that is difficult to refuse. \citet{graeber} showed that the moral language of debt and obligation is capable of expressing both exploitation and mutual care. \citet{polanyi} is often invoked on the side of Michéa, but his own position was that markets always operate within institutions, and that the question is which institutions are transformed or destroyed.

The synthesis that follows from this literature is that a viable society needs both impersonal protections and interpersonal obligations. Contracts and rights make commitments explicit and enforceable against parties with unequal power, while reciprocal practices sustain forms of cooperation that cannot be specified in advance. The defect of a market society, on Michéa's view, is the reduction of the second to the first. The defect of a purely communal society, on the evidence of the same literature, would be the reduction of the first to the second. Neither reduction is acceptable, and the essay's point is that the diagnosis of commodification requires attention to what is lost in a specific substitution and not a general preference for gift over contract.

\subsection{A formal remark on history and price}

The distinction between gift and contract has a structure that can be stated exactly. Let $\mathcal H_{\mathrm{social}}$ be the space of histories of interaction between agents, so that a history $h_{ab}=(e_1,\dots,e_n)$ records the events between two agents $a$ and $b$ in order. A market institution applies a projection
\[
\pi_{\mathrm{market}}:\mathcal H_{\mathrm{social}}\longrightarrow \mathcal X_{\mathrm{exchange}}
\]
into a space of exchangeable states, for instance the net price of services rendered. The projection is many-to-one. Two relationships with identical financial balances may differ radically in accumulated trust, resentment, gratitude, or obligation, and the projection maps them to the same point. Whatever is constitutive of the relationship beyond the balance is discarded.

This does not show that the projection is harmful. A projection is a compression, and compression is often the point of an institution. The harm arises when the compressed representation replaces the history as the object on which an institution acts, so that capacities that depended on the discarded information are no longer reproduced. Commodification, on this account, is the substitution of a projection for the history it summarizes, in settings where the discarded information was doing cooperative work. The formulation has the merit of being independent of any particular political valence. It applies equally to the commodification of tradition and of emancipation, and it makes the criterion for objection a matter of what the projection discards and whether the discarded capacities are regenerated elsewhere. The general form of this substitution is stated in Section~\ref{sec:mechanism}.

\section{Identity Labels and the Psychology of Classification}
\label{sec:labels}

The lecturer's most vivid claim concerns the commercialization of identity. Political and personal identities, he argues, are increasingly lived as sets of labels taken off the shelf, which commercial and political institutions address as market segments, and the person so classified becomes an ideal consumer of on-demand services. The comparison to choosing a member of a pop group is meant to show that the mechanism is the same.

Social psychology gives the claim a more careful footing. \citet{tajfel} showed that categorization alone, without material stakes, is sufficient to produce ingroup preference, and that individuals draw part of their self-concept from group membership. This means that labels are not trivial to the psychology of the persons who hold them. The theory also implies that the supply of salient categories shapes the identities available, and a commercial environment that continuously supplies segmented categories will shape the categories that become salient.

\citet{hacking86,hacking95} added a complication that works against a simple reading in either direction. Classifications of human beings are interactive: those classified become aware of the classification and alter their behavior in response, which in turn alters the classification. The looping effect means that labels cannot be treated as external impositions that distort an underlying authentic self, nor as neutral descriptions that leave the self untouched. A label can enable recognition, access, and coordination, and the same label can become a template that narrows the range of what the person is permitted to be. The essay therefore rejects the strong version of the lecturer's claim, on which minority recognition is reducible to advertising, and accepts a weaker version: commercial institutions can appropriate socially meaningful distinctions and convert them into instruments of differentiation, and doing so can shift the conditions under which identity becomes visible and actionable.

The theoretical significance is that segmentation is a case of the projection described above. Let $X$ be a person's situated history of relationships, commitments, and activities, and let $L$ be a finite space of labels, with an institutional classification $\lambda:X\to L$. The map is non-injective, so distinct persons receive identical classifications. The labels are useful for recognition and coordination. The difficulty appears when an institution treats $\lambda(x)$ as an exhaustive representation of $x$. At that point classification substitutes for relationship, and the person is legible to the institution through a set of attributes while the history that gives the attributes meaning is inaccessible. The same structure appears in credit scoring, employment screening, recommendation systems, and platform reputation, none of which concerns cultural politics. The mechanism is general, and it explains why the lecture's observation resonates beyond its stated subject. The general form is stated in Section~\ref{sec:mechanism}.

There is a puzzle in the lecture that this analysis helps to resolve. Pasolini, one of the preceding figures in the series, described the consumer society as a force of homologation that erased cultural difference and produced a uniform subject \citep{pasolini}. Michéa's thesis, as the lecturer presents it, is that consumer society produces a proliferation of identities. These appear to be opposite claims. They are reconcilable if the proliferation occurs at the level of labels and the homologation at the level of infrastructure, a distinction that is at least available, though whether it holds empirically is a further question. A society can be increasingly differentiated at the level of the segments it recognizes and increasingly standardized at the level of the systems that serve them. The distinction between surface differentiation and structural uniformity is, in this respect, more informative than either thesis taken alone.

Self-determination theory adds a last psychological consideration. \citet{deci} identified relatedness, alongside autonomy and competence, as a basic psychological need whose frustration impairs wellbeing. The thesis that dissolving durable relationships in favor of market-mediated substitutes frustrates relatedness is testable in principle, though the available evidence concerns specific interventions and not the large-scale claims made in the lecture. The essay notes the connection as an avenue of inquiry and does not treat it as confirmation.

\section{Literary Theory: Cash Nexus, Prefabricated Phrase, and the Plot of Emancipation}
\label{sec:literary}

The literary dimension of Michéa's lineage is more central than the lecture's emphasis on politics suggests, and it supplies a vocabulary for what the sociological and psychological arguments describe. \citet{carlyle} gave the Victorian critique of market society its enduring phrase when he insisted that payment in cash is not the sole nexus between persons, and the lecture invokes the term. The English tradition that runs from Carlyle through Ruskin and Morris to Orwell and, in a different register, Raymond \citet{williams}, treats the reduction of relations to payment as a failure of social description before it is a failure of policy. Literary criticism is here a form of social theory, because the question is which relations a language is capable of expressing.

Orwell's essay on the English language \citep{orwell46} provides a further connection that the lecture does not make. Orwell argued that prefabricated phrases and stale metaphors allow writers to dispense with thinking, because the phrase does the work that a perception should do. The lecturer's complaint about pre-made boxes, that the system supplies categories into which a person can be sorted and then ignored, has the same structure. A prefabricated label stands to the self as a stale metaphor stands to an observation: it saves effort at the cost of precision. This is not an argument against labels, which language cannot do without, but an argument about the conditions under which a label remains a tool and does not become a substitute for attention. The same objection can be directed at the lecture itself, whose own vocabulary of regimes, wokeness, and elites functions at times as a prefabricated phrase. A contested and shifting label such as the one the lecture employs is a weak instrument for analysis, and the structural claim defended here does not depend on it.

Pasolini's late journalism \citep{pasolini} extended the Orwellian insight into the claim that consumer society produced an anthropological mutation, a change in the human type that older forms of authoritarianism had failed to achieve. Whether or not the claim is accepted, it is a claim about the formation of subjects by an economic order, and it belongs with Michéa's thesis as one of its precursors. \citet{fisher2009} gave the contemporary form in his account of capitalist realism, the pervasive sense that no coherent alternative to capitalism can be imagined. Fisher is also a useful corrective to the lecture's appropriation of him, since his essay on the moralizing tendencies of the online left \citep{fisher2013} was a critique of left culture from within the left, directed at the replacement of solidarity by denunciation, and it did not endorse the claim that cultural liberation as such is a capitalist instrument.

The most interesting literary question concerns narrative form. The lecture reads a recent animated film as a parable in which a young woman is encouraged to abandon her family's tradition for a career in a large corporation, and presents the story as the embodiment of emancipation understood as exit. The reading of the film is contestable, and it is used here only as a statement of a narrative schema. The schema itself is worth analyzing. \citet{moretti} argued that the classical Bildungsroman was a symbolic form that reconciled two demands of modernity, the demand for individual self-determination and the demand for integration in a social order. The genre's central achievement was the representation of a socialization that did not feel like a loss of freedom. If this account is right, a narrative in which emancipation consists only of exit from the community has retained one half of the classical form and dropped the other. The sociological claim that the left has adopted a conception of freedom as release from constraint has a literary counterpart in the loss of the integrative half of the emancipation plot.

The related literary figure is the lesser evil. A politics organized as the avoidance of catastrophe has a characteristic plot, in which the present is justified by the horrors of the alternative, and the alternative is always located in the past or in a feared future. Fisher's analysis of capitalist realism and Michéa's analysis of the empire of the lesser evil converge on the observation that such a plot forecloses the question of what a good arrangement would be. In narrative terms, the lesser-evil plot has an antagonist and no protagonist, and political imagination atrophies in the absence of one.

\section{A Common Mechanism: Representation Substituted for Object}
\label{sec:mechanism}

The preceding sections have described three phenomena that were introduced separately. Orwell's prefabricated phrase substitutes for a perception, the identity label substitutes for a person, and the price substitutes for a relationship. In each case a representation of an object comes to stand in for the object in the operations that matter, and what the representation omits ceases to be available to those operations. The coincidence is not an accident of vocabulary, and it can be stated as a single structure.

Let $O$ be a domain of objects, which may be perceptions, persons, or relationships, and let $R$ be a space of representations with a representation map $\rho:O\to R$. An institution, which may be a writer's habit, a platform, or a market, is representation-bound when its action on an object depends on the object only through its representation, that is, when its decision rule has the form $g\circ\rho$ for some $g$ defined on $R$. Let $d$ be a distinction among objects that bears on some function $F$ on which the institution's operation relies: the accuracy of a sentence, the fit between a person and what is offered, the reliability of cooperation. The representation is lossy with respect to $d$ when objects that differ in $d$ receive the same representation. A representation-bound institution cannot respond to a distinction to which its representation is blind, and this is a structural fact, independent of anyone's intentions.

The second element is feedback. Objects in these domains are not inert. A person who is classified adjusts toward the classification, a relationship that is priced is redescribed as a service, and a writer's perceptions are shaped by the phrases available to express them. This can be written as a dynamic in which the next state of the object depends on its representation,
\[
o_{t+1}=f\bigl(o_t,\rho(o_t)\bigr),
\]
and it is the formal content of what Hacking called the looping effect of human kinds. Feedback alone does not entail that a distinction disappears. Whether the loss compounds depends on further conditions, which can be stated.

Let $\mu_d$ be the degree to which the distinction $d$ is maintained in $O$, and let $\mathcal I_d$ be the set of institutions that reward, record, or rehearse $d$, that is, institutions whose representations are sensitive to $d$. Suppose that maintenance requires support, so that $\mu_d$ decays at a rate $\lambda>0$ in the absence of it, that each institution $i\in\mathcal I_d$ contributes support $s_i\ge 0$ in proportion to the share of the relevant activity that it governs, and that the feedback of the preceding paragraph holds, so that objects adjust toward what the institutions governing them respond to. A simple form of these assumptions is
\[
\dot\mu_d=-\lambda\mu_d+\sum_{i\in\mathcal I_d}s_i .
\]

\begin{proposition}[Conditional erosion]\label{prop:erosion}
Under this dynamic $\mu_d$ converges to $\mu_d^{*}=\lambda^{-1}\sum_{i\in\mathcal I_d}s_i$. If the displacement of sensitive institutions by representation-bound institutions that are blind to $d$ reduces the total support below $\lambda\mu_d(0)$, then $\mu_d$ declines to a strictly lower level. If instead some institution in $\mathcal I_d$ retains support of at least $\lambda\mu_d(0)$, the decline is arrested.
\end{proposition}

\begin{proof}
The dynamic is linear, with solution $\mu_d(t)=\mu_d^{*}+\bigl(\mu_d(0)-\mu_d^{*}\bigr)e^{-\lambda t}$, and $\mu_d^{*}<\mu_d(0)$ if and only if $\sum_{i\in\mathcal I_d}s_i<\lambda\mu_d(0)$.
\end{proof}

The proposition is deliberately elementary. Its content is the identification of what the erosion claim depends on: that the distinction needs support, that the supporting institutions are displaced and not merely supplemented by blind ones, and that the objects respond to what is rewarded. Where an alternative institution preserves the distinction, the representation is lossy without being corrosive, and this is the role of regeneration channels elsewhere in the essay. This is also the way in which a classification can change the social reality it was introduced to describe. The claim of Michéa, as reconstructed here, is accordingly an empirical claim that, in specified domains, the supporting institutions have been displaced. It is not a consequence of the representation map alone.

The three cases instantiate the structure as follows. In Orwell's account, $O$ is the field of particular observations, $R$ is the stock of ready-made phrases, and the lost distinction is the specific feature of an observation that a fresh phrase would have been forced to register. In Hacking's account, $O$ is persons with their histories, $R$ is the set of available kinds, and the lost distinction is whatever in a life exceeds its category. In the market case, $O$ is relationships with their histories, $R$ is the space of exchangeable states, and the lost distinction is the accumulated history of reciprocity. Michéa's thesis, stripped of its polemical elaboration, is the claim that a growing share of cooperative activity is now mediated by institutions that are representation-bound in this sense, so that the history of the relationships through which cooperation was sustained is no longer registered, rewarded, or reproduced. The three cases are not metaphors for one another. They are instances of one structure, and this explains why the disciplines drawn on here converge. Sociology describes the substitution of price for reciprocity, psychology describes the substitution of label for self-understanding, and literary theory describes the substitution of phrase for perception, and each sees one instance and names it in its own vocabulary.

Not all substitution is harmful. Language cannot do without stock phrases, coordination among strangers requires labels, and prices transmit information that no individual could gather. Substitution becomes harmful under two conditions. The distinction $d$ that the representation discards must bear on a capacity on which the institution's own operation depends, and no other channel may regenerate that capacity. When both conditions hold, the institution consumes a stock it does not replenish, which is the sense of extraction used in this essay. Section~\ref{sec:criterion} states the condition as a criterion over histories.

\section{Politics: Class, Recognition, and the Question of a Criterion}
\label{sec:politics}

The political claim is that the left abandoned wages, rents, and ownership for cultural recognition, and so became a liberal pressure group whose demands are compatible with corporate capitalism. This is an empirical hypothesis about specific parties, institutions, and periods, and it cannot be established by observing that some corporations celebrate demographic representation. It can be examined through the distinction between changing who occupies positions within an institution and changing the institution's organization. A corporation can diversify its leadership without redistributing ownership or altering the relation between capital and labor.

The lecture tends to treat demographic inclusion and economic restructuring as competitors, but they are logically independent. \citet{fraserhonneth} argued that redistribution and recognition are two irreducible dimensions of justice, each of which can be denied without the other, and that political theory fails if it collapses either into the other. Access to housing, employment, education, and inheritance is shaped by both economic position and social discrimination, and a movement may address the two together. The legitimate institutional question in the lecture is whether organizations can adopt the language of social justice while leaving distributions of economic power unchanged. The answer depends on what changes are implemented and not on the vocabulary used to describe them.

\citet{lasch} made a related argument in the American context, suggesting that a professional and managerial class had detached itself from the common obligations of the communities that produced it, and had replaced a politics of shared institutions with a politics of mobility and self-presentation. The thesis resembles Michéa's, and it shares its principal weakness, which is that the opposition between a rooted populace and a mobile elite is better as a heuristic for noticing patterns than as a settled description of any population.

The politics of reception adds a final consideration. The lecture reports that Michéa has been cited by both the populist left and the populist right in France and has declined to denounce either. The position is intelligible as a refusal of guilt by association, and a philosopher is entitled to insist that his arguments be assessed on their merits and not on the identity of his readers. It also illustrates a general feature of the analysis above, which is that a diagnosis of social dissolution is available to be used by movements with opposed aims. The diagnosis does not select its own remedy, and the remedy depends on a criterion that the diagnosis does not supply. Section~\ref{sec:criterion} takes up that criterion.

\section{A Formal Statement of the Criterion of Viable Continuation}
\label{sec:criterion}

The criterion as first stated left three things unspecified: the time horizon over which viability is assessed, the meaning of viability itself, and whose cooperation must be preserved. A community can remain viable by excluding some of the people its rules govern, and a definition that leaves this open does not discriminate between common decency and well-organized domination. This section fixes the three parameters and builds exit and voice into the definition, instead of leaving them as separate normative qualifications.

\subsection{Setting}

Let $\mathcal A$ be the set of agents affected by a constraint structure $\mathcal C$, where an agent is affected if $\mathcal C$ governs its options or governs what others may do to it. The set is defined by effect and not by membership, so that people who are excluded from a community but subject to its rules belong to $\mathcal A$. Time is discrete, $t=0,1,\dots,T$, and unless otherwise stated $T$ is common knowledge among the agents, so that backward induction applies to finite-horizon interactions. A constraint structure fixes, for each agent, a set of permitted actions, together with the norms and sanctions that make some permitted actions more or less costly than others. A constraint is accordingly of one of two kinds: a restriction on the set of feasible actions, or a modification of the payoffs attached to feasible actions, such as a sanction. Removing a constraint means deleting a constraint of either kind. Let $\mathcal H(\mathcal C)$ be the set of joint histories $h=(s_0,\dots,s_T)$ that are sustained as equilibrium outcomes when each agent chooses among its permitted actions, in the sense that no agent gains by deviating given the others' behavior and the sanctions in force. Defining $\mathcal H(\mathcal C)$ by what is sustained, and not merely by what is permitted, is essential, because a constraint can enlarge the set of sustainable histories while restricting the set of permitted actions, and removing a constraint can do the opposite.

\subsection{Viability, exit, and voice}

Each agent $a$ has a capacity function $\kappa_a$ that assigns to each history prefix $h_{\le t}$ a real number measuring the agent's ability to sustain a recognizable course of life, with a threshold $\theta_a$ below which that ability fails. The community also has a stock $S$ of relational capacity, including trust, shared knowledge, and mutual assistance, with $S(h_{\le t})$ its level at time $t$. The set of viable histories at horizon $T$, where the admissible range of $T$ is fixed in the next subsection, is
\[
\mathrm{Viab}_T(\mathcal C)=\Bigl\{h\in\mathcal H(\mathcal C):\ \kappa_a(h_{\le t})\ge\theta_a\ \ \forall a\in\mathcal A,\ \forall t\le T;\ \ S(h_{\le t})\ge S(h_0)-\delta\ \ \forall t\le T;\ \ S(h_{\le\tau})\ge S(h_0)\ \ \forall\tau\in\mathcal T_T\Bigr\},
\]
where $\delta\ge 0$ is a tolerated depletion and $\mathcal T_T$ is the set of regeneration times not exceeding $T$, defined in the next subsection. The last two conditions together form the reproduction condition. The stock may fall within a regeneration cycle, by at most the tolerated amount $\delta$, but it must be restored to its initial level at each regeneration time, such as the end of a season or of a generation. A requirement that the stock never fall below its initial level would exclude systems that regenerate periodically, and a requirement that it be restored only at the final date would permit depletion that is merely deferred.

Let $\omega_a(h_{\le t})$ be the capacity agent $a$ would retain by leaving the arrangement at time $t$. A history is exit-secure if $\omega_a(h_{\le t})\ge\theta_a$ for every agent and every time, so that no affected agent is captive in the sense that leaving would push it below the threshold of viability. For each constraint $c$ in force at time $t$, let $f_c$ be the amendment rule that maps the stated positions of the agents it binds to a decision to retain or revise $c$. Responsiveness comes in two grades. The minimal grade requires only that there be some profile of the other agents' positions and two positions of $a$ that lead to different outcomes of $f_c$, which excludes dictatorship over $a$ but is compatible with a procedure that is inaccessible in practice. A history has voice in the sense used here only if it satisfies the effective grade. This grade is evaluated at the profiles that actually occur along the history. At each time, given the positions held by the agents outside some coalition containing $a$, that coalition must be decisive for $f_c$. All of its members must also be able to adopt $a$'s position at a cost that leaves their capacities above threshold. The effective grade implies the minimal grade, which is therefore a necessary condition. Neither grade requires that $a$ prevail. They require that its position can be made to count. Write $\mathrm{Exit}_T(\mathcal C)$ and $\mathrm{Voice}_T(\mathcal C)$ for the sets of histories in $\mathcal H(\mathcal C)$ that are exit-secure and that have voice at every time. The effective viable set is
\[
\mathrm{Viab}^{*}_T(\mathcal C)=\mathrm{Viab}_T(\mathcal C)\cap\mathrm{Exit}_T(\mathcal C)\cap\mathrm{Voice}_T(\mathcal C).
\]
Because exit and voice enter as conditions on membership in the set and not as additional commentary, a history in which cooperation is maintained by foreclosing the exit or the voice of some agent does not count as viable, however cohesive it is. This answers the objection that a community could remain viable while systematically excluding some participants. The excluded are in $\mathcal A$, and their capacities and options are subject to the same conditions as everyone else's.

\subsection{Horizon and comparison}

The horizon is not arbitrary. Let $T^{*}$ be the length of one regeneration cycle of the stock $S$, the time within which $S$ can be expected to be restored after depletion, and let $\mathcal T_T=\{kT^{*}:k\ge 1,\ kT^{*}\le T\}$ be the regeneration times not exceeding $T$. The criterion is stated for all $T\ge T^{*}$, and a claim about a transformation is credited only if it holds across that range. Two features keep this requirement from being stronger than intended. The reproduction condition is checked at regeneration times and not at every intermediate date, so that systems with periodic regeneration are not penalized for the troughs of their cycles, and the bounded-depletion condition with tolerance $\delta$ prevents the troughs from becoming unbounded. A transformation that enlarges the effective viable set at a short horizon and shrinks it once the horizon spans a regeneration cycle is extractive in the sense used above, because it borrows against $S$.

A transformation $\mathcal C\to\mathcal C'$ is an improvement at horizon $T$ if $\mathrm{Viab}^{*}_T(\mathcal C)\subseteq\mathrm{Viab}^{*}_T(\mathcal C')$, and a strict improvement if the inclusion is strict. All of these sets are subsets of a common universe of histories, so the comparison is well defined, and it is a partial order. Many transformations will be incomparable, since they enlarge the set in some directions and shrink it in others, and the criterion does not resolve such cases. It also does not derive the thresholds $\theta_a$ or the capacity functions $\kappa_a$, which are normative and empirical inputs and not outputs of the formalism. The criterion is therefore a discipline on argument. It requires that claims about the benefit or harm of removing a constraint be stated in terms of who is affected, over what horizon, and relative to what threshold, and it is not a procedure that settles political disputes.

\subsection{Two results}

\begin{proposition}[Non-monotonicity]\label{prop:nonmono}
There exist constraint structures $\mathcal C_1$ and $\mathcal C_2$ such that $\mathcal C_1$ is obtained from $\mathcal C_2$ by removing a constraint, so that every action feasible under $\mathcal C_2$ is feasible under $\mathcal C_1$ and the payoff modifications of $\mathcal C_1$ are a subset of those of $\mathcal C_2$, yet $\mathrm{Viab}^{*}_T(\mathcal C_1)\subsetneq\mathrm{Viab}^{*}_T(\mathcal C_2)$ for every $T\ge T^{*}$. The removed constraint can be either a restriction on feasible actions or a sanction.
\end{proposition}

\begin{proof}
Take two agents, $a$ and $b$, and a stage interaction in which in each period one of them, in alternation, suffers a shock that it cannot absorb alone, so that its capacity falls below its threshold unless it is assisted. The other agent may assist, at a cost $c>0$ to itself, which is smaller than the benefit $\beta>c$ to the agent assisted, and nothing else in the stage payoffs compensates the helper. Withholding therefore strictly dominates assisting in the stage game when no further constraint applies. Let $\mathcal C_1$ contain no constraint bearing on assistance. Over a known finite horizon the unique subgame-perfect equilibrium history is one in which assistance is never given, since withholding is dominant in the last period and the argument unravels backward through the earlier ones. The agent that suffers a shock then falls below its threshold, so $\mathrm{Viab}_T(\mathcal C_1)$ is empty. Let $\mathcal C_2$ add one constraint to $\mathcal C_1$, in either of two forms. In the first form it is a sanction, enforced at no cost to the agents, that lowers the payoff of withholding by an amount $s>c$, so that no feasible action is removed and only payoffs change. In the second form it is a prohibition that removes withholding from the feasible set. In both forms assisting is a best response in every period, the assisting histories are equilibria, neither agent falls below its threshold, and the stock of mutual assistance is maintained. Suppose also that both agents retain outside options above threshold and that the amendment rule over the added constraint is effectively responsive to each. Then these histories are exit-secure and have voice, so $\mathrm{Viab}^{*}_T(\mathcal C_2)$ is nonempty while $\mathrm{Viab}^{*}_T(\mathcal C_1)$ is empty. In the sanction form the feasible actions of the two structures coincide and $\mathcal C_1$ has strictly fewer payoff modifications, and in the prohibition form $\mathcal C_1$ has a strictly larger feasible set, so in each case $\mathcal C_1$ is obtained from $\mathcal C_2$ by removing a constraint.
\end{proof}

\begin{proposition}[Cohesion without exit or voice]\label{prop:cohesion}
If, for some history $h\in\mathcal H(\mathcal C)$, some agent $a\in\mathcal A$ is captive or lacks voice at some time $t\le T$, then $h\notin\mathrm{Viab}^{*}_T(\mathcal C)$, however high the capacities of the other agents and however large the stock $S$.
\end{proposition}

\begin{proof}
Immediate from the definition of $\mathrm{Viab}^{*}_T(\mathcal C)$ as an intersection that includes $\mathrm{Exit}_T(\mathcal C)$ and $\mathrm{Voice}_T(\mathcal C)$.
\end{proof}

\subsection{Consequences}

The criterion is symmetric between preservation and removal. An inherited constraint and a newly removed one are assessed by the same set, and the question in each case is the effect on the effective viable set of everyone affected. Proposition~\ref{prop:nonmono} shows why removal cannot be presumed to improve that set, and Proposition~\ref{prop:cohesion} shows why preservation cannot be presumed to. This is the sense in which the criterion distinguishes legitimate emancipation from destructive deregulation without presuming that inherited institutions are good or that newly recognized identities are commercial artifacts.

The criterion connects to the mechanism of Section~\ref{sec:mechanism} through the stock $S$ and the capacities $\kappa_a$. Where the conditions of Proposition~\ref{prop:erosion} hold for a distinction on which the regeneration of $S$ depends, an institution that is representation-bound with respect to that distinction will tend to produce histories that breach the depletion bound or fail the reproduction condition at a regeneration time, and such histories lie outside $\mathrm{Viab}_T$. This is the formal sense in which commodification is objectionable when it is, and the reason it is not objectionable merely because it is commodification. An activity is extractive when its return is obtained along histories that lie outside $\mathrm{Viab}_T$ because they deplete $S$. A platform that profits from communities it did not create, and optimizes in ways that deplete the relationships that attracted users, is extractive in this sense, whether the community in question is traditional or progressive. The criterion applies to the commercial channel that carries the lecture as it applies to anything else. A paid course can teach, a commercial publisher can circulate serious philosophy, and a monetized lecture can sustain a real intellectual community. The relevant question is whether the institution preserves the capacities underlying its relationships or systematically depletes them, and not whether it is commercial.

\section{Conclusion}

The lecture is a useful object of study because it presents a serious philosophical thesis in a form that mixes the thesis with its own polemical elaborations. Disentangled, the thesis yields four conclusions. The first is that the claim that capitalism can accommodate and commodify cultural emancipation is well supported, while the claim that emancipation is universally necessary to capitalism is not established, and the symmetry of the market's treatment of tradition and emancipation weighs against it. That symmetry leaves open historically specific causal relations between particular reforms and particular forms of market expansion, which require case-by-case evidence and are neither asserted nor rejected here. The second is that the central insight, available independently of the polemic, concerns representations that come to substitute for the objects they describe: the prefabricated phrase for the perception, the label for the person, the price for the relationship. The structure, with its feedback through which classifications reshape what they classify, is common to language, identity, and economic institutions, and it explains why sociology, psychology, and literary theory converge on the same phenomenon. The third is that Michéa's conservatism, which seeks to conserve the conditions of noncapitalist relationships, requires a criterion for distinguishing constraints that sustain reciprocal agency from constraints that preserve domination. The fourth is that such a criterion can be stated precisely. Defined over equilibrium histories, with an explicit horizon, a threshold of viability, a scope covering every affected agent, and exit and voice built in as conditions of membership, it is symmetric between preservation and removal, and it makes the position defensible without making it indiscriminate.

The essay does not establish that the doubling thesis is true, and it does not establish that it is false. It establishes that the thesis is better stated as a claim about reproduction and representation than as a claim about ideology, and that, so stated, it is compatible with a wide range of political commitments, including some that Michéa's most vocal readers would reject. The criterion offered here disciplines argument without settling it, and the erosion of a distinction under representation substitution is likewise conditional on the displacement of the institutions that sustain it. The thresholds of viability and the capacity functions are inputs that must be argued for, and many transformations will be incomparable under it. Freedom cannot be measured by counting the constraints that have been removed. It must also be evaluated in terms of the capacities, relationships, and future possibilities that remain available once they are gone, and that evaluation is the work that the doubling thesis, at its best, invites and does not complete.

\begin{thebibliography}{99}
\setlength{\itemsep}{0.4em}

\bibitem[{Academic Agent}(2026)]{academicagent}
Academic Agent (2026).
\newblock Why Consumerism Destroyed the Left: Jean-Claude Mich\'ea.
\newblock Video lecture, 8 October 2026. YouTube, channel Academic Agent.
\newblock \url{https://www.youtube.com/watch?v=xSUqEmBipL4}.

\bibitem[Bell(1976)]{bell}
Bell, D. (1976).
\newblock \emph{The Cultural Contradictions of Capitalism}.
\newblock New York: Basic Books.

\bibitem[{Boltanski and Chiapello}(2005)]{boltanski}
Boltanski, L., and Chiapello, \`E. (2005).
\newblock \emph{The New Spirit of Capitalism}.
\newblock Trans. G. Elliott. London: Verso.

\bibitem[Carlyle(1843)]{carlyle}
Carlyle, T. (1843).
\newblock \emph{Past and Present}.
\newblock London: Chapman and Hall.

\bibitem[{Deci and Ryan}(2000)]{deci}
Deci, E. L., and Ryan, R. M. (2000).
\newblock The ``What'' and ``Why'' of Goal Pursuits: Human Needs and the Self-Determination of Behavior.
\newblock \emph{Psychological Inquiry}, 11(4), 227--268.

\bibitem[Durkheim(1897)]{durkheim}
Durkheim, \'E. (1897).
\newblock \emph{Le Suicide: \'Etude de sociologie}.
\newblock Paris: Alcan.

\bibitem[Fisher(2009)]{fisher2009}
Fisher, M. (2009).
\newblock \emph{Capitalist Realism: Is There No Alternative?}
\newblock Winchester: Zero Books.

\bibitem[Fisher(2013)]{fisher2013}
Fisher, M. (2013).
\newblock Exiting the Vampire Castle.
\newblock \emph{openDemocracy}, 24 November 2013.

\bibitem[Fraser(2016)]{fraser2016}
Fraser, N. (2016).
\newblock Contradictions of Capital and Care.
\newblock \emph{New Left Review}, 100, 99--117.

\bibitem[Fraser(2017)]{fraser2017}
Fraser, N. (2017).
\newblock From Progressive Neoliberalism to Trump and Beyond.
\newblock \emph{American Affairs}, 1(4), 46--64.

\bibitem[{Fraser and Honneth}(2003)]{fraserhonneth}
Fraser, N., and Honneth, A. (2003).
\newblock \emph{Redistribution or Recognition? A Political-Philosophical Exchange}.
\newblock London: Verso.

\bibitem[Graeber(2011)]{graeber}
Graeber, D. (2011).
\newblock \emph{Debt: The First 5,000 Years}.
\newblock Brooklyn: Melville House.

\bibitem[Granovetter(1985)]{granovetter}
Granovetter, M. (1985).
\newblock Economic Action and Social Structure: The Problem of Embeddedness.
\newblock \emph{American Journal of Sociology}, 91(3), 481--510.

\bibitem[Hacking(1986)]{hacking86}
Hacking, I. (1986).
\newblock Making Up People.
\newblock In T. Heller, M. Sosna, and D. Wellbery (Eds.), \emph{Reconstructing Individualism}. Stanford: Stanford University Press.

\bibitem[Hacking(1995)]{hacking95}
Hacking, I. (1995).
\newblock The Looping Effects of Human Kinds.
\newblock In D. Sperber, D. Premack, and A. J. Premack (Eds.), \emph{Causal Cognition}. Oxford: Clarendon Press.

\bibitem[Haidt(2012)]{haidt}
Haidt, J. (2012).
\newblock \emph{The Righteous Mind: Why Good People Are Divided by Politics and Religion}.
\newblock New York: Pantheon.

\bibitem[Hirschman(1970)]{hirschman70}
Hirschman, A. O. (1970).
\newblock \emph{Exit, Voice, and Loyalty}.
\newblock Cambridge, MA: Harvard University Press.

\bibitem[Hirschman(1982)]{hirschman82}
Hirschman, A. O. (1982).
\newblock Rival Interpretations of Market Society: Civilizing, Destructive, or Feeble?
\newblock \emph{Journal of Economic Literature}, 20(4), 1463--1484.

\bibitem[Lasch(1995)]{lasch}
Lasch, C. (1995).
\newblock \emph{The Revolt of the Elites and the Betrayal of Democracy}.
\newblock New York: Norton.

\bibitem[Mauss(1990)]{mauss}
Mauss, M. (1990).
\newblock \emph{The Gift: The Form and Reason for Exchange in Archaic Societies}.
\newblock Trans. W. D. Halls. London: Routledge. (Original work published 1925.)

\bibitem[Mich\'ea(2007)]{michea}
Mich\'ea, J.-C. (2007).
\newblock \emph{L'Empire du moindre mal: Essai sur la civilisation lib\'erale}.
\newblock Paris: Climats.

\bibitem[Moretti(1987)]{moretti}
Moretti, F. (1987).
\newblock \emph{The Way of the World: The Bildungsroman in European Culture}.
\newblock London: Verso.

\bibitem[Orwell(1937)]{orwell37}
Orwell, G. (1937).
\newblock \emph{The Road to Wigan Pier}.
\newblock London: Victor Gollancz.

\bibitem[Orwell(1941)]{orwell41}
Orwell, G. (1941).
\newblock \emph{The Lion and the Unicorn: Socialism and the English Genius}.
\newblock London: Secker and Warburg.

\bibitem[Orwell(1946)]{orwell46}
Orwell, G. (1946).
\newblock Politics and the English Language.
\newblock \emph{Horizon}, 13(76), 252--265.

\bibitem[Pasolini(1975)]{pasolini}
Pasolini, P. P. (1975).
\newblock \emph{Scritti corsari}.
\newblock Milan: Garzanti.

\bibitem[Polanyi(1944)]{polanyi}
Polanyi, K. (1944).
\newblock \emph{The Great Transformation}.
\newblock New York: Farrar and Rinehart.

\bibitem[{Tajfel and Turner}(1979)]{tajfel}
Tajfel, H., and Turner, J. C. (1979).
\newblock An Integrative Theory of Intergroup Conflict.
\newblock In W. G. Austin and S. Worchel (Eds.), \emph{The Social Psychology of Intergroup Relations}. Monterey, CA: Brooks/Cole.

\bibitem[Thompson(1971)]{thompson}
Thompson, E. P. (1971).
\newblock The Moral Economy of the English Crowd in the Eighteenth Century.
\newblock \emph{Past and Present}, 50, 76--136.

\bibitem[Williams(1958)]{williams}
Williams, R. (1958).
\newblock \emph{Culture and Society, 1780--1950}.
\newblock London: Chatto and Windus.

\bibitem[Zelizer(2005)]{zelizer}
Zelizer, V. A. (2005).
\newblock \emph{The Purchase of Intimacy}.
\newblock Princeton: Princeton University Press.

\end{thebibliography}

\end{document}
